%%
%% This is file `sample-acmsmall.tex',
%% generated with the docstrip utility.
%%
%% The original source files were:
%%
%% samples.dtx  (with options: `all,journal,bibtex,acmsmall')
%% 
%% IMPORTANT NOTICE:
%% 
%% For the copyright see the source file.
%% 
%% Any modified versions of this file must be renamed
%% with new filenames distinct from sample-acmsmall.tex.
%% 
%% For distribution of the original source see the terms
%% for copying and modification in the file samples.dtx.
%% 
%% This generated file may be distributed as long as the
%% original source files, as listed above, are part of the
%% same distribution. (The sources need not necessarily be
%% in the same archive or directory.)
%%
%%
%% Commands for TeXCount
%TC:macro \cite [option:text,text]
%TC:macro \citep [option:text,text]
%TC:macro \citet [option:text,text]
%TC:envir table 0 1
%TC:envir table* 0 1
%TC:envir tabular [ignore] word
%TC:envir displaymath 0 word
%TC:envir math 0 word
%TC:envir comment 0 0
%%
%% The first command in your LaTeX source must be the \documentclass
%% command.
%%
%% For submission and review of your manuscript please change the
%% command to \documentclass[manuscript, screen, review]{acmart}.
%%
%% When submitting camera ready or to TAPS, please change the command
%% to \documentclass[sigconf]{acmart} or whichever template is required
%% for your publication.
%%
%%
%\documentclass[acmsmall,noacm]{acmart}
\documentclass[sigconf,nonacm]{acmart}
%\documentclass[11pt,a4paper,titlepage]{memoir}
%\documentclass{article}
%%
%% \BibTeX command to typeset BibTeX logo in the docs
\AtBeginDocument{%
  \providecommand\BibTeX{{%
    Bib\TeX}}}

% TODO: make a4

%% Rights management information.  This information is sent to you
%% when you complete the rights form.  These commands have SAMPLE
%% values in them; it is your responsibility as an author to replace
%% the commands and values with those provided to you when you
%% complete the rights form.
%\setcopyright{acmlicensed}
%\copyrightyear{2018}
%\acmYear{2018}
%\acmDOI{XXXXXXX.XXXXXXX}

%%
%% These commands are for a JOURNAL article.
%\acmJournal{JACM}
%\acmVolume{37}
%\acmNumber{4}
%\acmArticle{111}
%\acmMonth{8}
\usepackage{pgfplots}
\usepgfplotslibrary{statistics}
\pgfplotsset{compat=1.18}
\usepackage{subcaption}
\newcommand{\edgem}[2]{#1\rightarrow #2}
\newcommand{\edge}[2]{$\edgem{#1}{#2}$}
\newcommand{\wrt}{\textit{w.r.t.}}
\newcommand{\code}[1]{\texttt{#1}}
\newcommand{\codee}[1]{\texttt{#1}}
\newcommand{\ncite}{\footnote{missing citation}}
\usepgfplotslibrary{groupplots}
\usetikzlibrary{pgfplots.statistics}
\usepackage{booktabs}
\usepackage{multirow}
%\usepackage{siunitx}
\usepackage{pgfplotstable}

\newcommand{\recpart}{the Recursive Partitioning method}
\newcommand{\partitle}[1]{\textbf{#1:}\ }

\newcommand{\nitem}[1]{\item\partitle{#1}}


\geometry{a4paper}

\newcommand{\lorem}{

  Lorem ipsum dolor sit amet, consectetur adipisicing elit, sed do eiusmod
  tempor incididunt ut labore et dolore magna aliqua. Ut enim ad minim veniam,
  quis nostrud exercitation ullamco laboris nisi ut aliquip ex ea commodo
  consequat. Duis aute irure dolor in reprehenderit in voluptate velit esse
  cillum dolore eu fugiat nulla pariatur. Excepteur sint occaecat cupidatat non
  proident, sunt in culpa qui officia deserunt mollit anim id est laborum.

}

\usepackage{algorithm}
\usepackage{algpseudocode}

% Define functional-style keywords
\algnewcommand\algorithmicwhere{\textbf{where}}
\algnewcommand\Where{\algorithmicwhere}
\algnewcommand\algorithmiclet{\textbf{let}}
\algnewcommand\Let{\algorithmiclet}
\algnewcommand\algorithmiccase{\textbf{case}}
\algnewcommand\Case{\algorithmiccase}

\usepackage{listings}
\usepackage{xcolor}

\lstdefinelanguage{Haskell}{
  morekeywords={data, type, newtype, class, instance, where, let, in, case, of, if, then, else, do, module, import, deriving},
  sensitive=true,
  morecomment=[l]{--},
  morecomment=[s]{-}{-},
  morestring=[b]"
}
\DeclareMathOperator{\comps}{cc}
\DeclareMathOperator{\cost}{cost}
\DeclareMathOperator*{\argmin}{arg\,min}

\newcommand{\fc}{\mathit{fc}}
\newcommand{\mincost}{\operatorname{mincost}}


\usepackage[most]{tcolorbox}

\newtcolorbox{ghnote}{
    colback=blue!5,
    colframe=blue!60,
    boxrule=0pt,
    leftrule=3pt,
    arc=2pt,
    left=8pt, right=8pt, top=6pt, bottom=6pt,
    fontupper=\small,
}


\pgfplotsset{
  table/col sep=comma,
}

%%
%% Submission ID.
%% Use this when submitting an article to a sponsored event. You'll
%% receive a unique submission ID from the organizers
%% of the event, and this ID should be used as the parameter to this command.
%%\acmSubmissionID{123-A56-BU3}

%%
%% For managing citations, it is recommended to use bibliography
%% files in BibTeX format.
%%
%% You can then either use BibTeX with the ACM-Reference-Format style,
%% or BibLaTeX with the acmnumeric or acmauthoryear sytles, that include
%% support for advanced citation of software artefact from the
%% biblatex-software package, also separately available on CTAN.
%%
%% Look at the sample-*-biblatex.tex files for templates showcasing
%% the biblatex styles.
%%

%%
%% The majority of ACM publications use numbered citations and
%% references.  The command \citestyle{authoryear} switches to the
%% "author year" style.
%%
%% If you are preparing content for an event
%% sponsored by ACM SIGGRAPH, you must use the "author year" style of
%% citations and references.
%% Uncommenting
%% the next command will enable that style.
%%\citestyle{acmauthoryear}


%%
%% end of the preamble, start of the body of the document source.
\begin{document}

%%
%% The "title" command has an optional parameter,
%% allowing the author to define a "short title" to be used in page headers.
\title{Is Recursively Partitioned Search Space Enough for Operator Fusion?}

%%
%% The "author" command and its associated commands are used to define
%% the authors and their affiliations.
%% Of note is the shared affiliation of the first two authors, and the
%% "authornote" and "authornotemark" commands
%% used to denote shared contribution to the research.
\author{Marc Meier}
\affiliation{%
  \institution{ETH Zurich}
  \country{Switzerland}}
\email{marcme@student.ethz.ch}

%\author{Michel Weber}
%\authornote{Supervised by.}
%\email{michel.weber@inf.ethz.ch}
%\author{Zhendong Su}
%\authornotemark[1]
%\affiliation{%
%  \institution{ETH Zurich}
%  \country{Switzerland}}
%\email{zhendong.su@inf.ethz.ch}







%% Loading data
\pgfplotstableread[col sep=comma]{data/mnist/results.csv}          \dataRuntimeMnist
\pgfplotstableread[col sep=comma]{data/diffusion-large/results.csv}\dataRuntimeDiffusion
\pgfplotstableread[col sep=comma]{data/summary.csv}                \dataRecpart
\pgfplotstableread[col sep=comma]{data/summary-runtime.csv}        \dataRuntimeSummary
\pgfplotstableread[col sep=comma]{data/mnist/scaling.csv}          \dataScalingMnist
\pgfplotstableread[col sep=comma]{data/unet/scaling.csv}           \dataScalingUnet
\pgfplotstableread[col sep=comma]{data/cifar10/scaling.csv}        \dataScalingCifar
\pgfplotstableread[col sep=comma]{data/diffusion/scaling.csv}      \dataScalingDiffusion
\pgfplotstableread[col sep=comma]{data/diffusion-large/scaling.csv}\dataScalingDiffusionLarge
\pgfplotstableread[col sep=comma]{data/gpt/scaling.csv}            \dataScalingGpt
\pgfplotstableread[col sep=comma]{data/runtimes.csv}               \dataRuntimeSummary
\pgfplotstableread[col sep=comma]{data/relevance.csv}              \dataRuntimeRelevance



%%
%% By default, the full list of authors will be used in the page
%% headers. Often, this list is too long, and will overlap
%% other information printed in the page headers. This command allows
%% the author to define a more concise list
%% of authors' names for this purpose.
\renewcommand{\shortauthors}{Marc Meier}

%%
%% The abstract is a short summary of the work to be presented in the
%% article.
\begin{abstract}
  
  Operator fusion is a crucial optimization for computation graph compilers,
  as it optimizes kernel dispatches and memory operations to reduce
  the memory pressure of computation graph evaluators.
  Finding good operator fusions is an optimization problem with an exponential search space,
  which is usually solved via heuristics,
  which cannot, however, guarantee finding the global optimum due to their lack of exhaustiveness.
  We implemented a technique from the related problem of function inlining,
  designed to allow for exhaustive search for the global optimum by exploiting
  graph connectedness to reduce the search space.
  Our results indicated that this technique of recursively partitioned search space
  does not reduce the search space sufficiently on its own to make an exhaustive search feasible,
  as the reduced search space remains exponential with a comparatively large base.
  The search space remains large enough to require dropout ratios of up to 98\% on the largest models.
  We conclude that exhaustive search for optimal operator fusion requires a different partitioning strategy.

\end{abstract}

%%
%% The code below is generated by the tool at http://dl.acm.org/ccs.cfm.
%% Please copy and paste the code instead of the example below.
%%
%\begin{CCSXML}
%<ccs2012>
% <concept>
%  <concept_id>00000000.0000000.0000000</concept_id>
%  <concept_desc>Do Not Use This Code, Generate the Correct Terms for Your Paper</concept_desc>
%  <concept_significance>500</concept_significance>
% </concept>
% <concept>
%  <concept_id>00000000.00000000.00000000</concept_id>
%  <concept_desc>Do Not Use This Code, Generate the Correct Terms for Your Paper</concept_desc>
%  <concept_significance>300</concept_significance>
% </concept>
% <concept>
%  <concept_id>00000000.00000000.00000000</concept_id>
%  <concept_desc>Do Not Use This Code, Generate the Correct Terms for Your Paper</concept_desc>
%  <concept_significance>100</concept_significance>
% </concept>
% <concept>
%  <concept_id>00000000.00000000.00000000</concept_id>
%  <concept_desc>Do Not Use This Code, Generate the Correct Terms for Your Paper</concept_desc>
%  <concept_significance>100</concept_significance>
% </concept>
%</ccs2012>
%\end{CCSXML}

%\ccsdesc[500]{Do Not Use This Code~Generate the Correct Terms for Your Paper}
%\ccsdesc[300]{Do Not Use This Code~Generate the Correct Terms for Your Paper}
%\ccsdesc{Do Not Use This Code~Generate the Correct Terms for Your Paper}
%\ccsdesc[100]{Do Not Use This Code~Generate the Correct Terms for Your Paper}


%%
%% Keywords. The author(s) should pick words that accurately describe
%% the work being presented. Separate the keywords with commas.
%\keywords{Do, Not, Us, This, Code, Put, the, Correct, Terms, for,
%  Your, Paper}

%\received{20 February 2007}
%\received[revised]{12 March 2009}
%\received[accepted]{5 June 2009}


\begin{titlepage}
  \thispagestyle{empty}
  \centering
  \vspace*{1.5cm}
  {\scshape\large Bachelor's Thesis\par}
  \vspace{1.2cm}
  {\LARGE\bfseries Is Recursively Partitioned Search Space Enough for
   Operator Fusion?\par}
  \vspace{2.2cm}
  {\large Marc Meier\par}
  \vspace{0.35cm}
  {\large ETH Zurich, Switzerland\par}
  \vspace{1.8cm}
  {\itshape Supervised by\par}
  \vspace{0.35cm}
  {\large Michel Weber\par}
  {\large Prof.\ Dr.\ Zhendong Su\par}
  \vfill
  {\large \today\par}
\end{titlepage}

%%
%% This command processes the author and affiliation and title
%% information and builds the first part of the formatted document.
\maketitle

\section{Introduction}
  
  Operator fusion is a crucial optimization for computation graphs.
  Given a graph whose evaluation would dispatch many kernel launches,
  many of which do little more than loading data,
  performing a minor transformation on it,
  and storing it back to memory,
  it fuses them into a few kernels loading data once and
  doing many operators' computations on said data.
  
  Determining \emph{which} operators to fuse, however,
  is not trivial,
  as many fusions cause duplication of work,
  resulting in a balancing act,
  where too many fusions recompute the same value over and over again,
  while too few fusions waste performance on loading and storing large amounts
  of data from and to memory.

  As the total search space is infeasibly large,
  having $2^{|E|}$ where $|E|$ is the number of edges, possible fusion configurations,
  existing methods try to find good fusions using heuristics.
  However heuristics do not know the global optimum,
  as they are not exhaustive.
  A method exists for doing exhaustive search in a related problem (function inlining),
  but that method has not been applied to operator fusion.

  We applied that method, \recpart, to operator fusion.
  We introduced additional search space reductions,
  benchmarked resulting fusion configurations against the baseline,
  and modeled the scaling properties of the search space under \recpart.

  Our results indicate that the search space reduction of \recpart\ was not enough,
  even after adding additional reductions,
  as the search space still grows exponentially with a comparatively high base.

%\begin{tikzpicture}
%\pgfplotstableread[col sep=comma]{data/small/results.csv}\datatablesmall
%\begin{axis}[
%    boxplot/draw direction=y,
%]
% \addplot+[boxplot] table[y=baseline]{\datatablesmall};
% \addplot+[boxplot] table[y=ours]{\datatablesmall};
%\end{axis}
%\end{tikzpicture}

% Wie genau implementiert
% Related work
% Future work
% limitation
% discussion

\section{Background}

  This section provides the needed background on computation graphs,
  operator fusion, \recpart, and how a technique from function inlining
  translates to operator fusion.

  \subsection{Operator Fusion}

    This section explains computation graphs and operator fusion
    and the necessary terminology.

    \begin{figure}[t]%fig:compgraph
      \begin{subfigure}[b]{0.45\linewidth}
        \centering
        \begin{tikzpicture}
        \begin{scope}
          [ every node/.style={circle,draw,thick}
          , every edge/.style={draw,thick}
          ]
          \node (2)  at (0  ,1.5 ) {$2$};
          \node (A)  at (1  ,1.5 ) {$A$};
          \node (B)  at (2  ,1.5 ) {$B$};
          \node (t1) at (0.5,0.75) {$\times$};
          \node (t2) at (1.5,0.75) {$\times$};
          \node (p)  at (1  ,0   ) {$+$}[draw=blue];

          \path [->] (A) edge (t1);
          \path [->] (A) edge (t2);
          \path [->] (B) edge (t2);
          \path [->] (2) edge (t1);
          \path [->] (t1) edge (p);
          \path [->] (t2) edge (p);
        \end{scope}
        \end{tikzpicture}
        \caption{$2A+AB$}
        \label{fig:compgraph:expr}
      \end{subfigure}
      %
      \hfill
      %
      \begin{subfigure}[b]{0.45\linewidth}
        \centering
        \begin{tikzpicture}
        \begin{scope}
          [ every node/.style={circle,draw,thick}
          , every edge/.style={draw,thick}
          ]
          \node (A) at (0.5,1.5 ) {$A$};
          \node (B) at (1  ,0.75) {$B$};
          \node (T) at (0  ,0.75) {$\top$};
          \node (M) at (0.5,0   ) {$\bullet$}[draw=blue];

          \path [->] (A) edge (T);
          \path [->] (T) edge (M);
          \path [->] (B) edge (M);
        \end{scope}
        \end{tikzpicture}
        \caption{$A^\top B$}
        \label{fig:compgraph:linalg}
      \end{subfigure}


      \caption{Example: Showing how expressions can be represented as a computation graph.
      The marked node is the so-called \emph{root node} of the computation and represents the final output.}
      \label{fig:compgraph}
      \Description{Two computation graphs side by side, showing how expressions can be represented as computation graphs. The left represents $2A+AB$, the right represents $A^\top B$.}
    \end{figure}

    Computation graphs model computation as a directed graph,
    where nodes are operators transforming data and edges are the values
    flowing from one operator to the next
    \cite{survey2020},
    examples of which are shown in Figure~\ref{fig:compgraph}.
    As neural networks can naturally be expressed as computation graphs,
    where the operators are tensor operations working on large tensors,
    optimizing, compiling, and evaluating large computation graphs has become important.

    To leverage their high compute capabilities,
    tooling often evaluates such computation graphs on GPUs,
    by dispatching a compute kernel for every operator in the graph,
    passing it memory buffers storing their inputs,
    and memory buffers for them to store their outputs\cite{nvidiacudaguide}.
    This approach, however, results in each operator of the
    evaluated computation graph costing\cite{xla2023,survey2020}
    \begin{itemize}
      \item A memory read for the inputs
      \item A memory write for the result
      \item A kernel dispatch, with all the overhead that brings with it
    \end{itemize}

    For many of the cheaper operators, these costs \emph{greatly} overshadow
    the actual compute cost of the operator\cite{zhao2022apollo,tensorrt}.
    To address this, computation graph compilers employ an optimization called \emph{operator fusion},
    which fuses an operator \emph{producing} a value (the so-called \emph{producer})
    into an operator \emph{consuming} that value (the so-called \emph{consumer}),
    yielding a new operator, which performs the computation of \emph{both}
    unfused operators.
    This eliminates the need to store the result to memory, and to load it again for the consumer,
    saving the expensive memory interaction cost.

    For example, in Figure~\ref{fig:compgraph:linalg},
    the operator $\top$ (matrix transpose) can be fused into its only consumer,
    the matrix multiplication, yielding a fused operator, which performs matrix multiplication,
    % TODO: "transposed fashion" is weird wording
    but reads its left-hand-side in a transposed fashion, thus eliminating the need to
    generate and dispatch (i.e., \emph{materialize}) the transpose operator into its own kernel,
    saving both the kernel dispatch and the memory round-trip.

    Fusing an edge \edge{A}{B} on a computation graph
    leads to the following changes on the graph:
    \begin{itemize}
      \item
      % TODO: decide on a style
      $B$ is replaced with $B'$ (sometimes written $AB$),
      which is the new operator performing both $B$ \emph{and} $A$.

      \item
      The edge \edge{A}{B} is removed as the new operator
      $B'$ no longer needs the result of operator $A$ since it computes it directly.

      \item
      If after the above step $A$ no longer has any outgoing edges
      (i.e., $B$ was the only operator using $A$'s result),
      the node $A$ can be removed.
      Otherwise, $A$ remains since other operators still use its result.
      % TODO: we can only remove the vertex if it has no edges, not just no outgoing ones

      \item
      As the operator $B$ now computes what $A$ would have computed,
      $B$ inherits all dependencies on other operators that $A$ had,
      so the node $B'$ gains incoming edges from all direct predecessors of $A$.
    \end{itemize}


    \begin{figure}[t]%fig:fusing
      \begin{subfigure}[b]{0.45\linewidth}
        \centering
        \begin{tikzpicture}
        \begin{scope}
          [ every node/.style={circle,draw,thick}
          , every edge/.style={draw,thick}
          ]
          \node (A) at (0,1) {A};
          \node (B) at (1,0) {B};
          \node (C) at (2,1) {C};
          \node (D) at (3,0) {D};

          \path [->] (A) edge (B);
          \path [->] (B) edge (C)[draw=blue];
          \path [->] (B) edge (D);
          \path [->] (C) edge (D);
        \end{scope}
        \end{tikzpicture}
        \caption{Original}
        \label{fig:fusing:orig}
      \end{subfigure}
      %
      \hfill
      %
      \begin{subfigure}[b]{0.45\linewidth}
        \centering
        \begin{tikzpicture}
        \begin{scope}
          [ every node/.style={circle,draw,thick}
          , every edge/.style={draw,thick}
          ]
          \node (A) at (0,1) {A};
          \node (B) at (1,0) {B};
          \node (C) at (2,1) {C'};
          \node (D) at (3,0) {D};

          \path [->] (A) edge (B);
          \path [->] (A) edge (C);
          \path [->] (B) edge (D);
          \path [->] (C) edge (D);
        \end{scope}
        \end{tikzpicture}
        \caption{Fused}
        \label{fig:fusing:fused}
      \end{subfigure}

      \caption{Example: Showing the effect of fusing the edge \edge{B}{C}. The edge \edge{A}{B} is duplicated as the new edge \edge{A}{C'}. The newly fused $C'$ operator no longer depends on $B$, but it has inherited all dependencies of $B$, in this case $A$.}
      \label{fig:fusing}
      \Description{Two computation graphs side by side. The left shows the original graph, the right shows the graph after fusing the nodes B and C into BC}
    \end{figure}

    Partial fusions refer to fusions in which a producer is only fused into \emph{some} of its consumers,
    resulting in its kernel still getting materialized.
    An example of this is shown in Figure~\ref{fig:fusing},
    where $B$ only gets fused into $C$, but not $D$.
    The following comparison of the costs associated with $B$ before and
    after the partial fusion illustrate how a partial fusion can be beneficial,
    and why it often is not:
    \begin{itemize}
      \nitem{Before the partial fusion}
      \begin{itemize}
        \item Memory read for $B$'s inputs
        \item Compute cost to compute $B$'s result
        \item Memory write to store $B$'s result
        \item Memory read of $B$'s result for $C$
      \end{itemize}

      \nitem{After the partial fusion}
      \begin{itemize}
        \item Memory read for $B$'s inputs
        \item Compute cost to compute $B$'s result
        \item Memory write to store $B$'s result
        \item Memory read of $B$'s inputs for $C'$
        \item Compute cost to compute $B$'s result for $C'$
      \end{itemize}
    \end{itemize}
    All the costs with the materialized kernel of $B$ remain,
    but the cost of loading its result for $C$ is replaced
    with the cost of loading $B$'s inputs and computing its result inside $C'$.
    As a result, partial fusions improve performance only in the following scenarios:
    \begin{enumerate}
      \item
      When the combined cost of loading the inputs of $B$ and computing its result is lower than
      the cost of simply loading its result.

      \item
      When $C$ already depends on all of the inputs of $B$,
      the cost of loading its result only has to outweigh its compute cost.

    \end{enumerate}
    Notably, however, in case (1), the same argument can be made for
    fusing $B$ into \emph{any} of its consumers,
    so fusing $B$ into \emph{all} its consumers makes more sense.

    Whether a given edge can be fused is determined by the computation graph compiler's backend's
    ability to generate a performant kernel for a given fusion,
    so backend restrictions apply\cite{xla2023}.
    For example operators representing library functions from a vendor-optimized library
    cannot participate in operator fusion,
    as the backend cannot see into the kernel called by the library.

    Fusing all fusible edges, however, does not necessarily yield the best result,
    as fusions incur the following costs:
    \begin{itemize}
      \item
      As the kernels for fused operators are generated automatically by the backend,
      quality and performance of the generated kernels may be poor when fusing
      two complex operators.

      \item
      When fusing an operator into multiple consumers,
      the computation associated with the operator is effectively duplicated,
      as it is computed again by \emph{every} consumer.

    \end{itemize}
    Therefore, optimizing a computation graph entails finding good fusions,
    which presents an optimization problem.
    To tackle this optimization problem,
    we borrow a technique from function inlining as introduced in the next section.


  \subsection{Relating Function Inlining to Operator Fusion}\label{sect:correspondence}

    This section introduces the necessary function inlining concepts needed for \recpart,
    draws parallels between function inlining and operator fusion,
    and demonstrates how the graph modifications involved in the two transformations
    are equivalent.

    When compiling programs,
    a common optimization is \emph{function inlining},
    in which a function call site is replaced with the called function's body\cite{recpart-paper}.
    The goal is to
    \begin{itemize}
      \item
      remove function call overhead such as the called function's prologue and epilogue.

      \item
      extend the analysis scope of the compiler.
    \end{itemize}

    \begin{table*}[t]%tab:correspondence
      \caption{Showing correspondence of elements and operations from function inlining, operator fusion, and graph theory}
      \label{tab:correspondence}
      \begin{tabular}{ccc}
        \toprule
        Call graphs \& function inlining & Computation graphs \& operator fusion & Graph element or operation \\
        \midrule
        Function                      & Operator                                    & Vertex                                           \\
        Function call ($A$ calls $B$) & Data flow ($A$'s result is consumed by $B$) & Directed edge \edge{A}{B}                        \\
        $A$ is a caller of $B$        & $A$ is a producer to $B$                    & (\edge{A}{B}) $A$ is a direct predecessor of $B$ \\
        $B$ is called by $A$          & $B$ is a consumer of $A$                    & (\edge{A}{B}) $B$ is a direct successor of $A$   \\
        $B$ is inlined into $A$       & $A$ is fused into $B$                       & \edge{A}{B} merges                               \\
        \bottomrule
      \end{tabular}
    \end{table*}

    Function inlining operates on the program's \emph{call graph},
    which is a directed graph mapping each function to a vertex
    and each call site to a directed edge \edge{A}{B},
    where $A$ is the function \emph{containing} the call site,
    and $B$ is the function being \emph{called}.
    Therefore, inlining the call site \edge{A}{B} has the following effects on the call graph:
    \begin{itemize}
      \item
      $A$ is replaced with $A'$ (sometimes written $AB$),
      which is the body of $A$ with the call site of \edge{A}{B}
      replaced with the body of $B$.

      \item
      The edge \edge{A}{B} is removed as replacing the
      call site removed the call.

      \item
      If after the above step $B$ no longer has any incoming edges
      (i.e., \edge{A}{B} was the only call site calling $B$),
      the node $B$ can be removed.
      Otherwise, $B$ remains since other functions still call it.

      \item
      As the body of $B$ is inserted into the body of $A$,
      $A$ inherits all call sites residing in $B$,
      so the node $A'$ gains edges to all direct successors of $B$.
    \end{itemize}
    As the call graph modifications show,
    the operations performed on the call graph when inlining a function
    map one-to-one to the operations performed on a computation graph
    when fusing operators, but the direction of the directed edges is notably reversed.
    For example, inlining an edge \edge{A}{B} on a call graph performs \emph{the same} modification
    on the call graph as fusing an edge \edge{B}{A} would on a computation graph.
    Table~\ref{tab:correspondence} illustrates the correspondence.

  \subsection{Recursively Partitioned Search Space}

    When trying to find the optimal call sites to inline,
    a naïve approach is to evaluate every possible \emph{inlining configuration},
    which is an assignment of $\{\text{inline},\text{no-inline}\}$ to the call sites.
    This entails evaluating $2^n$ cases,
    where $n$ is the number of call sites,
    which, while feasible for smaller programs,
    becomes infeasible as the number of call sites grows.

    To reduce this search space,
    \citet{recpart-paper} introduced a technique that exploits the connectivity of
    the call graph to reduce the number of configurations evaluated,
    while remaining exact.

    Their key insight is that if a graph contains two or more connected components,
    each can be searched independently,
    as an inlining decision in one subgraph cannot affect an inlining decision in another.
    This considerably reduces the search space,
    since only $\sum_{g\in {cc}_G} |C_g|$ need to be evaluated,
    as opposed to the $\prod_{g\in {cc}_G} |C_g|$ a naive enumeration would visit,
    where $C_x$ are the searched inlining configurations of a graph $x$,
    and ${cc}_G$ are the connected components of $G$.

    Building on this insight, they proposed their technique of recursively partitioned search space
    (hereafter called \recpart),
    which models $C_G$ as a tree where each leaf represents an inlining configuration
    for the whole graph $G$ or a subgraph thereof.
    The tree is unfolded from the root by taking the graph $G$, picking an edge,
    and exploring both options recursively by modifying the graph accordingly:

    \begin{itemize}
      \nitem{$\text{inline}$}
      The edge is merged on the graph

      \nitem{$\text{no-inline}$}
      Not inlining the edge is identical to deleting it \wrt\ inlining,
      so the edge is deleted.
    \end{itemize}

    Both of the above mentioned modifications
    can split the graph into multiple connected components,
    specifically, if a bridge\footnote{
      A bridge is an edge of a graph,
      whose removal splits the graph into two connected components
    }
    was removed in the $\text{no-inline}$ case,
    or if a bridge \edge{A}{B} was merged where $B$ has no successors but other incoming edges.
    If this is the case, each connected component is recursively unfolded.

    Therefore, the tree consists of three kinds of nodes:
    \begin{itemize}
      \nitem{\code{InliningTreeLeafs}}
      A leaf where a decision has been made for every edge of the original graph
      or a subgraph thereof,
      representing an inlining configuration

      \nitem{\code{InliningTreeBinaryNodes}}
      Node with 2 children,
      representing exploring both cases for a given edge recursively

      \nitem{\code{InliningTreeComponentsNodes}}
      Node with multiple children,
      representing exploring all connected components of a disconnected graph independently.
    \end{itemize}

    It then folds that tree back into the optimal inlining configuration
    \wrt\ a cost metric to be minimized as follows:
    \begin{itemize}
      \nitem{\code{InliningTreeLeafs}}
      Evaluate the leaf inlining configuration on its cost,
      yielding both the cost and the configuration.

      \nitem{\code{InliningTreeBinaryNodes}}
      Compare the costs of both children,
      passing along the cost and configuration of the lower-cost one.

      \nitem{\code{InliningTreeComponentsNodes}}
      Combine the optimal configurations of all children into a collective one,
      evaluate its cost,
      and yield the combined configuration and the resulting cost
    \end{itemize}
    As in both steps, the unfolding and the folding, children of a given node
    are independent of each other, \recpart\ is embarrassingly parallel.

    While \citet{recpart-paper} uses binary size as a cost metric,
    the only requirement of the cost metric is
    \begin{equation}\label{equ:cost-req}
      \cost\left(G,\; \bigcup_{H \in {cc}_G} C^{*}_{H}\right)
      \;=\;
      \min_{C} \cost(G,\; C).
    \end{equation}
    where $\cost$ maps a graph and a configuration thereof to a the cost metric,
    $C^{*}_{g} \in \argmin_{C} \cost\left(g,\; C\right)$
    and ${cc}_G$ are the connected components of $G$.
    That is, \emph{optimality composes}:
    independently optimal configurations for the connected components of $G$
    form an optimal configuration for $G$.
    Note that the cost of the full graph does not necessarily need to be derivable
    from the costs of each connected component,
    so a separate evaluation of the combined cost is required,
    as is performed when folding an
    \code{InliningTreeComponentsNodes} node.

    As the search space reduction of this method is a direct result of exploring
    connected components separately,
    it is important to pick edges whose removal produces multiple connected components often.
    Thus, care must be taken picking a good bridge if a graph has bridges
    and picking a good edge in the absence of bridges.
    For this purpose, they employed heuristics to pick central bridges,
    and pick edges that reduce high out degrees and create as many bridges as possible
    if no bridge exists.

    As a result, speedup heavily depends on graph topology, as that determines the frequency of bridges.
    For example in a simple line
    $A\rightarrow B\rightarrow C\rightarrow D\rightarrow E\rightarrow F$
    every single edge is a bridge, so on every iteration it gets to split the graph.
    A star like
    $\{\edgem{A}{B},\edgem{A}{C},\edgem{A}{D},\edgem{A}{E},\edgem{A}{F}\}$
    on the other hand,
    would see no speedup as this graph would never contain a bridge.

    Using the recursive partitioning method,
    they showed that they could reduce the search space for optimally inlining
    SPEC2017 from $2^{349}$ to $2^{25}$,
    enabling an exhaustive search.

    Because of the correspondence pointed out in Section~\ref{sect:correspondence},
    the recursive partitioning method can be applied to find the optimal
    operator fusion configuration of computation graphs.



\section{Methods}

  This section explains
  (1) how the XLA computation graph compiler was modified to
  facilitate controlling and evaluating fusions,
  (2) how the search itself was implemented,
  (3) what graphs were used to test the search on,
  (4) how the results of the search were validated against real measurements,
  and (5) how the scaling properties of the graphs were measured.

  \subsection{XLA}

    XLA is a computation graph compiler and evaluator
    that can be targeted by tensorflow, JAX and PyTorch
    and supports running models on CPUs, GPUs and other accelerators.
    It consumes so-called HLO modules,
    which can contain multiple computation graphs,
    with one being marked as the root computation of the module.
    Each computation is defined using their HLO operations,
    connected using static single assignment (SSA),
    some of which can reference other computations,
    such as reductions over tensors and more importantly fusions.

    Fusions are described by referencing another computation in the same module
    and having the backends convert that other computation to a single operator,
    in the GPU backend's case turning that computation into a singular kernel.

    Like many other compilers,
    XLA compiles computation graphs via passes analyzing or modifying a module,
    before passing it to a target-specific backend for final lowering.
    While many passes of the pipeline generate fusion instructions,
    most of the fusions are generated in a fusion pipeline. % TODO: verify

    To generate fusions ourselves
    and to facilitate interacting with XLA in general,
    we created 3 passes running at key steps in XLA's pipelines.
    Each pass takes some arguments from the process environment,
    returning immediately when its parameters are not defined in the environment
    without performing any modifications to the module,
    effectively disabling the pass:
    \begin{itemize}
      \nitem{\code{ForwardFusion}} dumps the state of the module very early on before many optimizations
      to a file specified in its environment-passed argument

      \nitem{\code{ForceFusion}} reads a list of fusions to perform from an
      environment-passed filename,
      performing those fusions in order.
      It runs just after \code{ForwardFusion}.

      \nitem{\code{QualityMetric}} uses XLA's internal runtime estimation mechanisms
      to estimate the runtime of the root computation in nanoseconds,
      returning that value through an environment-passed filename.
      XLA's internal runtime estimation mechanisms take a description of a GPU
      to model runtime on that GPU\cite{openxlagpumodel}.
      It runs after most optimizations,
      at the latest point before many operations are rewritten using
      GPU specific intrinsics.
    \end{itemize}

    As fusions are such a crucial optimization for computation graphs,
    many different passes, even normalization passes, generate fusions.
    Instead of preventing fusions which we don't force,
    we let the compiler produce fusions after our forced fusions
    as it normally would.
    This has the helpful side-effect,
    that passing an empty list of fusions is equivalent to
    not having changed anything about fusions,
    serving as a baseline for XLA's fusion process.

  \subsection{Implementation of \recpart}

    The search for the optimal fusion configuration is implemented
    via an external harness executing the XLA utility \code{hlo\_opt} to
    \begin{itemize}
      \item dump the graph to perform \recpart\ on.
      \item evaluate configurations by forcing its fusions and having it
      estimate the runtime with the given fusions.
    \end{itemize}
    The harness implements \recpart\ with the following
    notable differences from \citet{recpart-paper}:
    \begin{itemize}
      \item
      Instead of separately unfolding the tree and then evaluating it,
      it implements one recursion doing both.

      \item
      Edge selection differs slightly in that when a graph contains multiple bridges,
      it picks the one whose deletion produces the most balanced split
      between the two subgraphs in terms of edge count.

      \item
      Instead of minimizing a cost metric,
      the nanosecond runtime $c$, as returned from XLA,
      is converted to a \emph{quality} metric $q$
      as $q=\frac{1}{c}$,
      so as to allow returning a quality metric of 0 in cases where \code{hlo\_opt} fails,
      which can happen when nonsensical fusions are explored.
      As \recpart\ only considers comparisons of different cost metrics
      and negative cost metrics cannot occur,
      this change does not change the optimization as long as
      the comparisons are inverted.

      \item
      When propagating up the most optimal configuration,
      if both candidate configurations have the same quality metric,
      the one with fewer forced fusions is propagated up.
      This is done such that the configuration forcing no fusions,
      which is, as previously stated, equivalent to the baseline,
      is the final returned optimal configuration if no configuration
      manages to improve upon it,
      making it easy to dismiss results that fail to improve on the baseline.
      Note that this has no effect on the final quality metric,
      as it only changes what configuration is returned if the other candidate has
      \emph{the same} quality metric.
    \end{itemize}

    As discussed in Section~\ref{sect:discussion},
    \recpart\ is not enough to reduce the search space to a feasible degree,
    so the following additional optimizations are introduced to further cut down
    the search space:
    \begin{itemize}
      \item
      XLA defines some operators as not being able to participate in fusions.
      All edges to and from such operators are removed before \recpart.

      \item
      An option called \code{fuse-all-consumers} is introduced to only consider fusions where operators
      are fused into all their consumers,
      thus partial fusions are not considered.
      The same reduction is used by Optimus\cite{cai2022optimus}. % TODO: is this a good idea?
      It does not change how edges are selected,
      instead it just merges or splits all edges that share
      the producer of the selected edge.

      \item\label{sect:dropout}
      A \emph{dropout} parameter is introduced to prune a specified ratio of the graph.
      Edges to prune are selected by sorting the edges topologically,
      dropping edges from both the front and back of this edge list until enough edges are dropped.
      If the resultant graph has multiple connected components,
      the largest one in terms of edge count is searched by \recpart.

      % TODO: maybe this belongs somewhere else?
      While this has ramifications for the validity of the results
      as explained in Section~\ref{sect:limitations},
      the size of the search space searched with dropout is a lower bound
      for the size of the full search space,
      as every configuration in the set of searched fusion configurations with dropout
      maps to at least one distinct configuration in the set of searched fusion configurations without dropout.
    \end{itemize}

    The searches were run on an AMD Ryzen Threadripper 3990X with up to 32 threads in parallel.

  \subsection{The modules that were optimized}
    The HLO modules searched for optimality were generated by
    implementing common neural network architectures with Keras
    and dumping their evaluation calls.
    The architectures and their models used are as follows:
    \begin{itemize}
      \nitem{Image classifier}
      Two sizes of image classifier were used:
      A small MNIST handwritten digit classifier and a larger Cifar10 classifier

      \nitem{Image segmentation}
      A UNet image segmentation architecture

      \nitem{Image Diffusion}
      Two sizes of image diffusion models using a UNet architecture
      with transformer blocks in their bottlenecks.

      \nitem{Generative Pretrained Transformer (GPT)} % TODO: do this for all items with titles
      Three sizes of autoregressive GPT models,
      differing in dimensionality of latent token representation
      and differing dimensionality of the Key- and Value-vector
      and the number of heads in their attention blocks.
      As they only differ in the dimensionality of the tensors,
      the graphs are equivalent,
      but different fusion configurations might still be optimal
      as costs and benefits of different fusions may fall out differently.
    \end{itemize}
    All the models used are comparatively small,
    the largest of which, Diffusion (L), is just under 10 Million parameters.
    Bigger models failed to fit into the RTX 2080 Ti's memory.

    \begin{table}[t]%tab:models
      \caption{Summary of the models optimized, accompanied by the dropout ratio used}
      \label{tab:models}
      \centering
      \pgfplotstabletypeset[
        col sep=comma,
        columns={name, pre_edges, dropout},
        %
        columns/dropout/.style={
          column name=Dropout,
          column type=r,
          preproc/expr={100*##1},
          postproc cell content/.append style={
            /pgfplots/table/@cell content/.add={}{\%},
          },
          fixed, fixed zerofill, precision=1,
        },
        %
        columns/name/.style      ={column name=Model            , string type, column type=l},
        columns/pre_edges/.style ={column name=Edges        , int detect, column type=r},
        %
        every head row/.style={
          before row={%
            \toprule
          },
          after row=\midrule,
        },
        every last row/.style={after row=\bottomrule},
      ]{\dataRecpart}
    \end{table}

    For each model,
    a dropout value was picked by using a \emph{dry run} mode of the harness
    where it estimates search time,
    targeting a total search time of just over 10 hours.
    Search times were estimated by counting the number of evaluations,
    using an estimated evaluation rate of 60 evaluations per second.
    The picked dropout values are listed in Table~\ref{tab:models},
    even just marginally reducing these dropout values quickly
    grows estimated search time to multiple days.
    All models were searched with the \code{fuse-all-consumers} option enabled.

  \subsection{Benchmarking}
    After the harness returned optimal fusions to force for each model,
    the ones for which it reported an improvement over the baseline
    according to its quality metric were benchmarked against the baseline
    on an Nvidia RTX 2080 Ti using XLA's \code{run\_hlo\_module} utility,
    timing 200 runs after 25 warmup runs of both the baseline and the optimized module.

  \subsection{Scaling}
    To quantify how the number of configurations to evaluate scales in terms of
    the number of edges in the graph,
    sweeps of the dropout rate were run for the models\footnote{
      As the graphs for the three GPT models are identical and the number
      of evaluated configurations does not depend on the evaluation result,
      the quality metric,
      the number of evaluated configurations for a given dropout rate
      are identical for all three GPT models
      and as a result, a sweep of one of them applies to all three.
    },
    recording both edge count and number of evaluated configurations.


\section{Results and Discussion}\label{sect:discussion}

  \begin{table*}[t]%tab:results
    \caption{Results for each model optimized}
    \label{tab:results}
    \centering
    %\resizebox{\textwidth}{!}{%
      \pgfplotstabletypeset[
        col sep=comma,
        columns={name, dropout, pre_edges, post_edges, evals, r_eval, b_metric, o_metric, i_metric},
        %
        columns/dropout/.style={
          column name={},
          column type=r,
          preproc/expr={100*##1},
          postproc cell content/.append style={
            /pgfplots/table/@cell content/.add={}{\%},
          },
          fixed, fixed zerofill, precision=1,
        },
        columns/i_metric/.style={
          column name=Impr.,
          column type=r,
          preproc/expr={100*##1},
          fixed, fixed zerofill, precision=3, skip 0.,
        },
        columns/r_eval/.style={
          column name={},
          column type=c,
          postproc cell content/.append style={
            /pgfplots/table/@cell content/.add={}{\;e/s},
          },
          fixed, fixed zerofill, precision=1, skip 0.,
        },
        postproc cell content/.append code={%
          \ifnum\pgfplotstablecol=8\relax
            \if\relax\detokenize{#1}\relax\else
              \pgfkeys{/pgfplots/table/@cell content/.add={}{\%}}%
            \fi
          \fi
          \if\relax\detokenize{#1}\relax
            \pgfkeyssetvalue{/pgfplots/table/@cell content}{\textemdash}%
          \fi
        },
        %
        columns/name/.style      ={column name={}             , string type, column type=l},
        columns/pre_edges/.style ={column name=Initial        , int detect, column type=r},
        columns/post_edges/.style={column name=After          , int detect, column type=r},
        columns/evals/.style     ={column name={}             , int detect, column type=r},
        columns/b_metric/.style  ={column name=Baseline       , sci, sci zerofill, sci precision=4, column type=r},
        columns/o_metric/.style  ={column name=Optimal        , sci, sci zerofill, sci precision=4, column type=r},
        %
        every head row/.style={
          before row={%
            \toprule
            \multirow{2}{*}{Model} &
            \multirow{2}{*}{Dropout} &
            \multicolumn{2}{c}{Edges} &
            \multirow{2}{*}{\shortstack{Evaluated\\configs}} &
            %\multirow{2}{*}{\shortstack{Evaluations\\per second}} &
            \multirow{2}{*}{\shortstack{Eval\\rate}} &
            \multicolumn{3}{c}{Quality metric} \\
            \cmidrule(lr){3-4} \cmidrule(lr){7-9}
          },
          after row=\midrule,
        },
        every last row/.style={after row=\bottomrule},
      ]{\dataRecpart}
    %}
  \end{table*}

  Table~\ref{tab:results} shows results of the harness runs.
  The \emph{Edges} columns show the number of edges in the graph both before and after dropout.
  The columns under \emph{Quality metric} contain the quality metrics of
  both the baseline and the optimized module as computed by the harness,
  showing the relative improvement on the quality metric in cases where
  an improvement could be found.
  The search yielded marginal improvements for only two models.
  The \emph{Eval rate} column lists the average evaluation rate for a given model,
  yielding a total search time of just under 35 hours for all models.
  The drop from the expected evaluation rate of 60 e/s down to
  11.4 e/s in the worst case can be explained by the work not being split up evenly
  between the execution threads.

  \begin{table}[t]%tab:runtime
    \caption{Runtimes of the two models benchmarked}
    \label{tab:runtime}
    \centering
    %\resizebox{\textwidth}{!}{%
      \pgfplotstabletypeset[
        col sep=comma,
        columns={stat_name,m_baseline,m_optimal,m_impr,d_baseline,d_optimal,d_impr},
        %
        columns/m_impr/.style={
          column name=Impr.,
          column type=r,
          preproc/expr={100*(##1-1)},
          fixed, fixed zerofill, precision=3, skip 0.,
        },
        columns/d_impr/.style={
          column name=Impr.,
          column type=r,
          preproc/expr={100*(##1-1)},
          fixed, fixed zerofill, precision=3, skip 0.,
        },
        postproc cell content/.append code={%
          \ifnum\pgfplotstablecol=3\relax\if\relax\detokenize{#1}\relax\else
              \pgfkeys{/pgfplots/table/@cell content/.add={}{\%}}%
          \fi\fi
          \ifnum\pgfplotstablecol=6\relax\if\relax\detokenize{#1}\relax\else
              \pgfkeys{/pgfplots/table/@cell content/.add={}{\%}}%
          \fi\fi
          \if\relax\detokenize{#1}\relax
            \pgfkeyssetvalue{/pgfplots/table/@cell content}{--}%
          \fi
        },
        %
        columns/stat_name/.style={column name={}, string type, column type=l},
        columns/m_baseline/.style={column name=Basel., fixed, fixed zerofill, precision=4, skip 0., column type=r},
        columns/m_optimal/.style ={column name=Opt.  , fixed, fixed zerofill, precision=4, skip 0., column type=r},
        columns/d_baseline/.style={column name=Basel., fixed, fixed zerofill, precision=4, skip 0., column type=r},
        columns/d_optimal/.style ={column name=Opt.  , fixed, fixed zerofill, precision=4, skip 0., column type=r},
        %
        every head row/.style={
          before row={%
            \toprule
             &
            \multicolumn{3}{c}{MNIST} &
            \multicolumn{3}{c}{Diffusion (L)} \\
            \cmidrule(lr){2-4} \cmidrule(lr){5-7}
          },
          after row=\midrule,
        },
        every last row/.style={after row=\bottomrule},
      ]{\dataRuntimeSummary}
    %}
  \end{table}

  \begin{figure}[t]%fig:runtime
    \begin{tikzpicture}
    \begin{groupplot}[
      group style={
        group size=2 by 1,
        horizontal sep=0,
      },
      boxplot/draw direction=y,
      boxplot/whisker range=100,
      width=0.5\linewidth, height=0.75\linewidth,
      xtick={1,2},
      xticklabels={Baseline, Optimal},
      boxplot/.append style={
        draw=black,
        %thick,
      }
    ]

      \nextgroupplot[
        title=MNIST,
        axis y line*=left,
        %yticklabel pos=left,
      ]
      \addplot+[boxplot, boxplot/draw position=1] table[y=baseline]{\dataRuntimeMnist};
      \addplot+[boxplot, boxplot/draw position=2] table[y=ours]{\dataRuntimeMnist};

      \nextgroupplot[
        title=Diffusion (L),
        %axis y label=right,
        yticklabel pos=right,
        ytick pos=right,
      ]
      \addplot+[boxplot, boxplot/draw position=1] table[y=baseline]{\dataRuntimeDiffusion};
      \addplot+[boxplot, boxplot/draw position=2] table[y=ours]{\dataRuntimeDiffusion};

    \end{groupplot}
    \end{tikzpicture}


    \caption{
      The distribution of measured runtimes of both
      the \emph{MNIST} and \emph{Diffusion (L)} models in seconds.
    }
    \label{fig:runtime}
    \Description{
      A box-and-whisker plot showing measured runtimes of both
      the \emph{MNIST} and \emph{Diffusion (L)} models in seconds.
    }
  \end{figure}

  Because the harness predicted improvements for the MNIST and Diffusion (L) models,
  these two were benchmarked for real runtime performance,
  yielding runtimes as shown in Figure~\ref{fig:runtime},
  the mean, median and standard deviation of which are listed in Table~\ref{tab:runtime}.
  Figure~\ref{fig:runtime} shows that the MNIST runtime exhibits rather long tails,
  while Diffusion (L) has shorter tails.
  It is notable that both models benchmarked had rather short runtimes,
  both being under 0.25 seconds,
  likely contributing to the high degree of noise present in both distributions,
  as overheads related to dispatching likely play a relatively large role in the
  runtimes.

  \begin{table}[t]%tab:relevance
    \caption{P-Value and 95\% Confidence Interval of Median Improvement}
    \label{tab:relevance}
    \centering
    %\resizebox{\textwidth}{!}{%
      \pgfplotstabletypeset[
        col sep=comma,
        columns={benchmark,p_val,micil,miciu},
        %
        columns/miciu/.style={
          column name=Upper,
          column type=r,
          preproc/expr={100*(##1-1)},
          fixed, fixed zerofill, precision=1, skip 0.,
          postproc cell content/.append style={
            /pgfplots/table/@cell content/.add={}{\%},
          },
        },
        columns/micil/.style={
          column name=Lower,
          column type=r,
          preproc/expr={100*(##1-1)},
          fixed, fixed zerofill, precision=1, skip 0.,
          postproc cell content/.append style={
            /pgfplots/table/@cell content/.add={}{\%},
          },
        },
        %
        columns/benchmark/.style={column name={}, string type, column type=l},
        columns/p_val/.style={column name={}, fixed, fixed zerofill, precision=2, skip 0., column type=r},
        %
        every head row/.style={
          before row={%
            \toprule
             &
            \multirow{2}{*}{P-Value} &
            \multicolumn{2}{c}{Median Impr. CI} \\
            \cmidrule(lr){3-4}
          },
          after row=\midrule,
        },
        every last row/.style={after row=\bottomrule},
      ]{\dataRuntimeRelevance}
    %}
  \end{table}

  Evaluating the statistical significance of the runtime results using a
  Mann-Whitney U test and a 95\% confidence interval for improvement in terms
  of median runtime, as seen in Table~\ref{tab:relevance},
  it becomes apparent that the differences visible in Table~\ref{tab:runtime}
  are not statistically significant.
  With P-Values greater than the conventional .05 limit,
  a statistically significant difference in the runtime distributions could be determined
  for neither of the benchmarked models.
  This is also supported by the 95\% confidence interval for median runtime improvement
  including 0\% for both benchmarks.

  Comparing predicted improvement based on the quality metric as shown in Table~\ref{tab:results}
  to measured runtimes,
  we observe that the predicted improvement of .749\% for the MNIST model
  failed to materialize in the measured runtime,
  pointing to the evaluated quality metric being an inaccurate proxy
  for real runtime.
  Notably, the same claim cannot be made for the Diffusion (L) model,
  as its predicted improvement of .015\% is within the 95\%
  confidence interval for that model, as shown in Table~\ref{tab:relevance}.

  \begin{figure*}[t]%fig:scaling
    \newcommand{\scalingplot}[4]{%
      \begin{subfigure}[b]{#4}
        \centering
        \begin{tikzpicture}
        \begin{axis}[
          width=\linewidth,
          ymode=log,
          log basis y=2,
          xmin=1,
          ymin=1,
          grid=both,
        ]
          \addplot[only marks, mark size=1px] table[x=edges, y=evals                                   ]{#3};
          \addplot[no markers               ] table[x=edges, y={create col/linear regression={y=evals}}]{#3};
        \end{axis}
        \end{tikzpicture}
        %
        \pgfmathparse{2^\pgfplotstableregressiona}
        \xdef\regressionbase{\pgfmathresult}
        \pgfmathparse{2^\pgfplotstableregressionb}
        \xdef\regressionconstant{\pgfmathresult}
        %
        \caption{%
          \centering
          #1 \\
          Regression parameters $cb^x$:\\
          ($c=\regressionconstant$\;$b=\regressionbase$)%
        }
        \label{#2}
      \end{subfigure}%
    }
    \scalingplot{MNIST        }{fig:scaling:mnist}          {\dataScalingMnist         }{0.32\linewidth}
    \scalingplot{UNet         }{fig:scaling:unet}           {\dataScalingUnet          }{0.32\linewidth}
    \scalingplot{Cifar10      }{fig:scaling:cifar}          {\dataScalingCifar         }{0.32\linewidth}
    \scalingplot{Diffusion (S)}{fig:scaling:diffusion-small}{\dataScalingDiffusion     }{0.32\linewidth}
    \scalingplot{Diffusion (L)}{fig:scaling:diffusion-large}{\dataScalingDiffusionLarge}{0.32\linewidth}
    \scalingplot{GPT (All)    }{fig:scaling:gpt}            {\dataScalingGpt           }{0.32\linewidth}

    \caption{
      Plotting edges after dropout against evaluated configs on a log-linear plot with a linear regression fitted for each model
    }
    \label{fig:scaling}
    \Description{
      Six log-linear scatter plots plotting the relation between edge count and search space size for the six different
      computation graphs.
      They all appear to be mostly linear or piecewise linear with both vertical and horizontal jumps, indicating exponential growth with differing regimes, each of which has differing growth rates.
    }
  \end{figure*}

  Figure~\ref{fig:scaling} plots the six sweeps run on the models on a log-linear plot,
  each accompanied by a linear regression finding parameters $c$ and $b$
  of the exponential function $cb^x$ where $x$ is the number of edges.
  The bases $b$ are in the range from $1.2$ to $1.6$.

  \subsection{Search Space is still too Large}
    While \recpart\ did reduce the search space,
    it did not yield as big of a gain as \citet{recpart-paper}.
    Even with the additional search space reduction of not considering partial fusions,
    a high dropout ratio had to be used to get results in feasible time.
    It appears that the search space remains infeasible due to the following factors reinforcing one another:
    \begin{itemize}
      \nitem{Large graphs}
      % TODO: "classical programming" is shakey
      As compared to classical programming where functions consist of multiple instructions,
      expressing computations as computation graphs tends to need many more operators
      connected together as compared to functions calling each other in classical programming,
      because functions tend to express a lot more complex logic as they already are a composition
      of multiple fundamental instructions.
      While computation graphs can also express more complex operations using a single operator,
      such as fusion operations being a combination of the operators they fuse,
      many of these \emph{composite} operators appear as a result of optimization passes,
      so they tend to be scarce in unoptimized graphs.
      % TODO: XLA spams operators maybe?

      \nitem{Search space still grows exponentially}
      As most log-linear plots from Figure~\ref{fig:scaling} appear linear,
      the search space seems to still grow exponentially \wrt\ the number of edges in the graph,
      and unfortunately,
      looking at the regressed bases $b$, while below two,
      are still large enough,
      that the size of the search space explodes,
      even for relatively small graphs.
    \end{itemize}


  \subsection{Scaling}

    This subsection discusses patterns visible in the plots in
    Figure~\ref{fig:scaling},
    reasoning about what causes them
    and what they might tell us.

    Either horizontal or vertical jumps can be observed in most of the models.
    Horizontal jumps can be explained by how the dropout mechanisms reduces the graph.
    As explained in Section~\ref{sect:dropout},
    dropout sorts the edges topologically,
    removes edges from both ends of the sorted edge list,
    and picks the biggest connected component of the remaining graph.
    This has as effect that there can be a second, smaller,
    connected component that is pruned by the dropout logic.
    If a reduced dropout rate now reveals a bridge between two connected components,
    the two previously disconnected connected components merge into one,
    instantly increasing the edge count of the searched graph by both the
    bridge connecting them and the entirety of the smaller connected component
    which was previously pruned.
    Considering, however, the edge connecting the newly connected components is a bridge,
    \recpart\ likely picks that bridge as its first edge,
    instantly reducing the search space again,
    explaining why horizontal jumps do not seem to also jump vertically.

    Vertical jumps may be explained by quirks in the edge selection heuristic
    rather than graph structure,
    as vertical jumps seem to curve away from the jump,
    sometimes both before and after the jump.
    Figure~\ref{fig:scaling:unet} for example shows,
    that the vertical jump above 30 edges is immediately followed by a
    \emph{course correction} back to the regressed line.
    Similarly, Figure~\ref{fig:scaling:cifar} exhibits a dip just before the vertical jump,
    leading to the vertical jump \emph{itself} being the course correction back
    to the regressed line.
    The jump just below 50 edges in Figure~\ref{fig:scaling:mnist}
    exhibits both the dip before the vertical jump,
    as well as the course correction back down after the jump.


\section{Limitations}\label{sect:limitations}
  
  % TODO: in the right place insert ramifications / limitations of using cost metric as proxy for performance

  While not invalidating findings about search space size,
  the high dropout ratios have ramifications on findings in the following ways:
  \begin{itemize}
    \item
    As dropout prunes edges, it naturally has an impact on graph topology which,
    as previously discussed, % TODO: actually discuss this "previously"
    can change how aggressively \recpart\ reduces the search space.
    As a result, extrapolating the scaling statistics to scenarios
    without dropout is not possible

    \item
    Since dropout reduces the graph visible to \recpart,
    it may hide potential regions which \emph{could} be optimized and which
    \emph{can} improve on the baseline,
    so claims about the lack of improvements cannot be made.
  \end{itemize}

  % TODO: Maybe a paragraph about that there are _multiple_ fusion strategies to choose for each fusion
  %       which was not explored?

  Similarly,
  the decision to only optimize using the reduced search space of not considering partial fusions % TODO: guess who's back
  may hide potential more optimal solutions, again, invalidating claims about lack of improvement. % TODO: ditto
  % TODO: maybe add something about edge selection in fuseall mode
  % TODO: maybe say something about the limited number of models benchmarked


\section{Related work}

  This section categorizes related work into:
  (1) differing strategies for finding fusions,
  (2) approaches for reducing search spaces for compiler optimizations,
  and (3) how exhaustive searches are enabled for other domains.

  \partitle{Operator Fusion Strategies in Compilers}
  Operator Fusion is established as a crucial optimization for computation graph compilers
  \cite{survey2020},
  seeing different strategies in use,
  such as rule-based approaches that apply a fixed set of transformations whenever their
  preconditions match\cite{chen2018tvm, xla2023},
  pattern-matching approaches finding common specific patterns of
  well-known operator pairs and replacing them with equivalent operators
  \cite{tensorrt},
  exposing scheduling decisions to the user\cite{ragan2013halide},
  to more modern approaches specialized to fill specific gaps of the previously mentioned
  approaches\cite{niu2021dnnfusion, zhao2022apollo, chimera2023, zhao2026neptune}.
  These strategies iteratively address specific failure modes of existing solutions,
  but none of these attempt to find the globally optimal fusion configuration.
  They are designed to be fast and predictable rather than exhaustive.
  This gap motivates our investigation.

  \partitle{Search Space Reduction for Compiler Optimizations}
  Several efforts exist that try to reduce the search space to feasible sizes via heuristic methods,
  such as using machine learning models\cite{ashouri2018survey, cui2025decos} and
  pruning candidates by collapsing multiple into groups\cite{phothilimthana2016lens}.
  While these efforts find good solutions fast, they do so by pruning unlikely candidates,
  making them non-exhaustive and incapable of finding the global optimum with certainty.

  \partitle{Exhaustive search and \recpart} % TODO: find better title
  Exhaustive search guarantees optimality, but only under conditions that keep the search space bounded.
  \citet{recpart-paper} exhaustively search optimal inlining decisions by
  exploiting call graph connectivity to reduce the space from $2^{349}$ to $2^{25}$.
  HieraSynth\cite{lu2025hierasynth} achieves completeness in super-optimization by
  hierarchically partitioning candidate programs and pruning unrealizable branches,
  making exhaustive search feasible for instruction sets up to 700 instructions.
  Bik\cite{bik2021exhaustive} decouples search time from compile time by performing
  exhaustive state space search during compiler development,
  and Optimus\cite{cai2022optimus} achieves exhaustive search by both
  requiring a cost metric that decomposes additively over the chosen partition
  and exploiting a maximum out-degree that is small in practice.
  These approaches share a common pattern: exhaustiveness is bought by restricting the problem
  to a bounded graph structure,
  a fixed cost model, or an offline search.
  Whether they hold for general computation graphs under a runtime-based objective,
  rather than only for restricted graph structures and static cost models,
  is the question our work investigates.
  

\section{Future work}

  % TODO: feels clunky
  This section lists gaps in our results and how they might be expanded into
  future research.

  \partitle{Investigate search space scaling}
  Our results indicate that the search space scales a lot worse when compared to \citet{recpart-paper}.
  Future work could investigate what makes computation graphs scale worse than call graphs
  by observing \emph{how} \recpart\ carves up the graph and how topologies found in
  computation graphs affect its ability to reduce the search space

  \partitle{Non-partial fusion specific edge selection}
  Our edge-selection heuristic was intended to find good edges when only considering the
  fusion of that specific edge.
  Future work could attempt to find an edge-selection heuristic better suited for
  searching with the \code{fuse-all-consumers} option enabled,
  by considering that fusing or not fusing a given edge also applies that decision to
  all outgoing edges of the producer.

  \partitle{More representative metric}
  We observed that using a model as a search metric, as a proxy for a different metric
  can produce non-representative search results.
  A more direct metric, such as benchmarked runtime, could be used to find improvements
  that a model does not model.

  \partitle{Additional reduction in search space}
  We adapted a search space reduction strategy from call graphs which can contain cycles.
  As computation graphs cannot contain cycles,
  future work could find a way to exploit this for an additional reduction in search space.

\section{Conclusion}
  
  We applied \recpart\ to operator fusion in XLA with the goal of exhaustively finding
  optimal fusion configurations.
  Our results indicate, that the search space reduction provided by \recpart\ is not sufficient,
  even after additional search space reductions,
  to exhaustively search optimal fusion configurations.
  On the heavily pruned graphs
  we could search, we found no statistically significant improvement.
  This is not evidence that no improvement exists,
  given the exclusion of partial fusions, the heavy pruning, and a
  demonstrably inaccurate quality metric.
  As a result we conclude that \recpart\ alone is not sufficient for operator fusion
  on the computation graphs we tested.
  The method's feasibility for inlining depends on a static and accurate cost metric
  and a call graph structure that partitions well.
  Neither seems to be true for computation graphs,
  where runtime performance statically can only be estimated
  and search spaces still grow quickly even after application of \recpart.

%%
%% The acknowledgments section is defined using the "acks" environment
%% (and NOT an unnumbered section). This ensures the proper
%% identification of the section in the article metadata, and the
%% consistent spelling of the heading.
%\begin{acks}
%  To Robert, for the bagels and explaining CMYK and color spaces.
%\end{acks}

%%
%% The next two lines define the bibliography style to be used, and
%% the bibliography file.
\bibliographystyle{ACM-Reference-Format}
\bibliography{bib}


%%
%% If your work has an appendix, this is the place to put it.
\appendix

\section{Repositories}
  The XLA fork is found at https://github.com/0x000C0A71/xla,
  the harness source code at https://github.com/0x000C0A71/RecursivePartitioning,
  the scripts generating the models at https://github.com/0x000C0A71/ThesisModels,
  and the thesis sources including the raw results at https://github.com/0x000C0A71/Thesis

%\section{Runtime data of the individual runs}\label{sect:raw-runtime}
%  \pgfplotstabletypeset[
%    col sep=comma,
%    create on use/mnist_base/.style={copy column from table={\dataRuntimeMnist}{baseline}},
%    create on use/mnist_ours/.style={copy column from table={\dataRuntimeMnist}{ours}},
%    create on use/diff_base/.style={copy column from table={\dataRuntimeDiffusion}{baseline}},
%    create on use/diff_ours/.style={copy column from table={\dataRuntimeDiffusion}{ours}},
%    columns={mnist_base, mnist_ours, diff_base, diff_ours},
%    % Column types and formatting
%    columns/mnist_base/.style={column name={Baseline}, skip 0., fixed, fixed zerofill, precision=6, column type=r},
%    columns/mnist_ours/.style={column name={Ours    }, skip 0., fixed, fixed zerofill, precision=6, column type=r},
%    columns/diff_base/.style ={column name={Baseline}, skip 0., fixed, fixed zerofill, precision=6, column type=r},
%    columns/diff_ours/.style ={column name={Ours    }, skip 0., fixed, fixed zerofill, precision=6, column type=r},
%    % Inject the two-row grouped header
%    every head row/.style={
%        before row={%
%            \toprule
%            \multicolumn{2}{c}{MNIST} &
%            \multicolumn{2}{c}{Diffusion} \\
%            \cmidrule(lr){1-2} \cmidrule(lr){3-4}
%        },
%        after row=\midrule,
%    },
%    every last row/.style={after row=\bottomrule},
%  ]{\dataRuntimeMnist}

%\section{Research Methods}
%  \subsection{Part One}
%    \lorem
%
%  \subsection{Part Two}
%    \lorem
%
%  \section{Online Resources}
%    \lorem

\end{document}
\endinput
%%
%% End of file `sample-acmsmall.tex'.
